% Generated draft: requires mathematical and source review.
\documentclass[11pt]{article}
\usepackage[margin=25mm]{geometry}
\usepackage{fontspec}
\setmainfont{Latin Modern Roman}
\usepackage{amsmath,amssymb,mathtools,unicode-math}
\setmathfont{Latin Modern Math}
\usepackage{enumitem,needspace,xurl,hyperref}
\setlength{\parindent}{0pt}
\setlength{\parskip}{5pt}
\setlength{\emergencystretch}{3em}
\newcommand{\N}{\mathbb{N}}
\newcommand{\Z}{\mathbb{Z}}
\newcommand{\Q}{\mathbb{Q}}
\newcommand{\R}{\mathbb{R}}
\newcommand{\C}{\mathbb{C}}
\newcommand{\F}{\mathbb{F}}
\newcommand{\A}{\mathbb{A}}
\newcommand{\PP}{\mathbb{P}}
\newcommand{\E}{\mathbb{E}}
\newcommand{\eps}{\varepsilon}
\newcommand{\dd}{\,\mathrm{d}}
\begin{document}
\section*{UnsolvedMath ID: 30001493}
\textbf{Source problem code:} catalogue entry
\subsection{Definitions and mathematical statement}
Let \(M\) be a complex manifold. A Kähler form \(\omega\) is a closed, positive real \((1,1)\)-form on \(M\). In the usual unnormalized convention, a Kähler--Ricci flow is a smooth family of Kähler forms satisfying
\[
\frac{\partial\omega(t)}{\partial t}=-\operatorname{Ric}(\omega(t)),
\qquad
\operatorname{Ric}(\omega)=-\sqrt{-1}\,\partial\bar\partial\log\det(g_{j\bar k}),
\]
where \((g_{j\bar k})\) is the associated Hermitian metric in local holomorphic coordinates. This equation fixes terminology only: the supplied record does not specify whether the intended flow is normalized or unnormalized.

\textbf{Source question.} For a Kähler--Ricci flow with singular time \(T\), determine its behavior as \(t\uparrow T\). In particular, determine whether, and under what hypotheses, the flow can pass through the singularity and continue on a new manifold.

\textbf{Unresolved scope.} The source does not specify the initial manifold or metric, compactness or projectivity assumptions, whether \(T\) is finite, the definition of singular time, or the topology and regularity in which limiting behavior is sought. It also does not define passage through a singularity: the continuation category, the relation between the original and subsequent spaces, and the compatibility condition at \(T\) are unspecified. In particular, it is unclear whether the subsequent space must be a smooth complex manifold or may be a singular variety. Thus this is a research question with unresolved hypotheses and continuation criteria, rather than a fully quantified existence assertion.
\subsection{Short English statement}
Describe how a Kähler--Ricci flow behaves near a singular time and determine when it can continue through the singularity on a new space. The supplied wording leaves the geometric hypotheses and the meaning of continuation unspecified. Its supplied status is partially solved; this draft makes no independent literature assessment.
\subsection{Sources}
\begin{itemize}
\item Supplied local UnsolvedMath registry record, permanent ID \(30001493\), source problem code \texttt{catalogue entry}, title: \emph{Continuation of Kähler--Ricci Flow through Singularities}. The record belongs to the Oberwolfach Reports Open Problems collection and gives the proposed year as 2010; no individual proposer is supplied.
\item \emph{Geometrie} (2010), Oberwolfach Reports, pp.\ 1535--1572, DOI 10.4171/owr/2010/26. Supplied source report and literature reference, labeled ``original OWR Report 26/2010'': \url{https://doi.org/10.4171/owr/2010/26}. The supplied background describes these reports as documenting mathematical workshops and their open-problem discussions.
\item Supplied literature reference, labeled ``Song--Tian weak Kähler--Ricci flow through singularities'': \url{https://arxiv.org/abs/0909.4898}.
\item Literature assessment dated exactly \texttt{2026-08-21}: the supplied curation labels the problem \texttt{partially\_solved} and states that Song--Tian constructed unique weak Kähler--Ricci flows through singularities on projective varieties with log-terminal singularities, while the broad continuation question in the extracted wording remains unsettled in full generality. This is supplied dated curation, not an independently verified assessment. The supplied background explicitly requests independent verification of time-sensitive status and cited claims.
\item The registry labels the statement \texttt{corrected\_verified} and reports that it was checked against the official Oberwolfach report, with the clean statement preserving the stated hypotheses, notation, and open target. This verification claim is supplied provenance; no independent source consultation is asserted here.
\end{itemize}
\end{document}
